\section{Fermat varieties and their Hodge characters}\label{sec:fermat}

\subsection{Correspondences}
All varieties are smooth and projective over $\CC$, possibly disconnected,
and cohomology has rational coefficients unless we say otherwise. Let $B$
and $X$ be varieties, with $B$ of pure dimension $b$. An algebraic cycle
$\Gamma$ of codimension $c$ on $B\times X$ gives the map
\[
  \Gamma_{*}\colon H^{k}(B)\to H^{k+2c-2b}(X),\qquad
  y\mapsto\pr_{X*}\bigl(\cl(\Gamma)\cup\pr_{B}^{*}y\bigr),
\]
a morphism of Hodge structures $H^{k}(B)(b-c)\to H^{k+2c-2b}(X)$. It carries
algebraic classes to algebraic classes, the composite of two such maps is the
map of the composite correspondence, and for a morphism $f\colon B\to X$ the
graph of $f$ and its transpose give $f_{*}$ and $f^{*}$
\cite[Chapter~16]{Ful98}.

A pure rational Hodge structure $V$ of weight $w$ is \emph{polarizable} if
it carries a bilinear form $Q$ with $Q(Cx,\bar x)>0$ for $x\ne0$, where $C$
acts on $V^{a,b}$ by $i^{a-b}$. The cohomology of a variety and its Tate
twists are polarizable \cite[Chapter~6]{Voi02b}, and polarizable Hodge
structures are semisimple.\footnote{If $W\subseteq V$ is a sub-Hodge
structure, the restriction of $Q$ to $W$ is again a polarization, so it is
nondegenerate and $V=W\oplus W^{\perp}$. The orthogonal $W^{\perp}$ is a
sub-Hodge structure, because $V_{\CC}=\bigoplus V^{a,b}$ is an orthogonal
decomposition for the hermitian form $Q(Cx,\bar y)$, and $W_{\CC}$ is the sum
of its intersections with the $V^{a,b}$.}

\begin{lemma}\label{lem:lift}
Let $\varphi\colon V\to W$ be a surjective morphism of polarizable Hodge
structures of weight $2p$. Then $\varphi$ maps the Hodge classes of $V$, the
rational classes of type $(p,p)$, onto those of $W$.
\end{lemma}

\begin{proof}
By semisimplicity the kernel of $\varphi$ has a complement $V'$ that is a
sub-Hodge structure, and $\varphi$ restricts to an isomorphism $V'\to W$ of
Hodge structures. Its inverse carries a Hodge class of $W$ to a Hodge class
of $V'\subseteq V$.
\end{proof}

\begin{theorem}[Lefschetz \cite{Lef24}, Kodaira--Spencer
\cite{KS53}]\label{thm:lef11}
Every Hodge class of degree two is algebraic.
\end{theorem}

\begin{proof}
In the exact sequence $H^{1}(X,\cO_{X}^{\times})\to H^{2}(X,\ZZ)\to
H^{2}(X,\cO_{X})$ of the exponential sequence, the second map is the
composite of $H^{2}(X,\ZZ)\to H^{2}(X,\CC)$ with the projection to
$H^{0,2}(X)$. So an integral class of type $(1,1)$ is the first Chern class
of a holomorphic line bundle, which is algebraic by GAGA, and its Chern class
is the class of a divisor. A rational Hodge class has an integral multiple.
\end{proof}

\subsection{The Fermat varieties}
For $m\ge1$ and $n\ge0$ let $X^{n}_{m}\subseteq\PP^{n+1}$ be the Fermat
variety $x_{0}^{m}+\dots+x_{n+1}^{m}=0$; thus $X^{0}_{m}$ is a set of $m$
points and $X^{1}_{m}$ the Fermat curve. The group
$G_{n}=\mu_{m}^{n+2}/\mu_{m}$ acts by multiplying the coordinates, and its
characters are the vectors $a=(a_{0},\dots,a_{n+1})\in(\ZZ/m)^{n+2}$ with
$\sum_{i}a_{i}=0$. Let $\mathfrak A^{n}_{m}$ be the set of characters with
every $a_{i}\ne0$. For $y\in\ZZ/m$ let $\bar y\in\{0,\dots,m-1\}$ be its
representative and $\langle y\rangle=\bar y/m$, and for
$a\in\mathfrak A^{n}_{m}$ put $\langle a\rangle=\sum_{i}\langle a_{i}\rangle$,
an integer between $1$ and $n+1$.

\begin{proposition}\label{prop:fermatchars}
\begin{enumerate}[label=\textup{(\roman*)}]
\item The primitive cohomology $H^{n}_{\mathrm{prim}}(X^{n}_{m},\CC)$ is the
direct sum of lines $V(a)$, $a\in\mathfrak A^{n}_{m}$, on which $G_{n}$ acts
by pull-back through the character $a$, and $V(a)\subseteq H^{n-q,q}$ with
$q=\langle a\rangle-1$.
\item Let $n$ be even. The space of primitive Hodge classes, tensored with
$\CC$, is the sum of the lines $V(a)$ with $\langle ta\rangle=n/2+1$ for
every $t\in(\ZZ/m)^{\times}$.
\end{enumerate}
\end{proposition}

\begin{proof}
(i) By Griffiths' description of the cohomology of a hypersurface by
residues \cite[\S8]{Gri69}, $H^{n-q,q}_{\mathrm{prim}}(X^{n}_{m})$ is
isomorphic to the graded piece of degree $(q+1)m-(n+2)$ of the Jacobian ring
$\CC[x_{0},\dots,x_{n+1}]/(x_{0}^{m-1},\dots,x_{n+1}^{m-1})$; the monomial
$x^{b}$ corresponds to the residue of $x^{b}\Omega/F^{q+1}$, where
$F=\sum x_{i}^{m}$ and $\Omega=\sum_{i}(-1)^{i}x_{i}\,dx_{0}\wedge\dots
\widehat{dx_{i}}\dots\wedge dx_{n+1}$. That residue is an eigenvector of
$G_{n}$ with character $a=b+(1,\dots,1)$, as $\Omega$ has character
$(1,\dots,1)$ and $F$ is invariant. As $b$ runs over the exponents with
$0\le b_{i}\le m-2$, $a$ runs once over the characters with all entries
nonzero, and the degree condition reads $\langle a\rangle=q+1$.

(ii) The decomposition is defined over $\QQ(\mu_{m})$, and the automorphism
$\zeta\mapsto\zeta^{t}$ of $\QQ(\mu_{m})$ carries $V(a)$ to $V(ta)$. So for
each orbit $O$ of $(\ZZ/m)^{\times}$ on $\mathfrak A^{n}_{m}$ the sum
$\bigoplus_{a\in O}V(a)$ is defined over $\QQ$, and these rational subspaces
decompose $H^{n}_{\mathrm{prim}}(X^{n}_{m},\QQ)$. A rational class with a
nonzero component on some $V(a)$ has nonzero components on the whole orbit of
$a$; so it is of type $(n/2,n/2)$ only if every $V(ta)$ is, that is, only if
$\langle ta\rangle=n/2+1$ for every unit $t$.
\end{proof}

\subsection{The inductive structure}
Fix $\varepsilon$ with $\varepsilon^{m}=-1$. For $r,s\ge1$, let
$Y\subseteq X^{r}_{m}\times X^{s}_{m}$ be the subvariety $x_{r+1}=y_{s+1}=0$,
which is isomorphic to $X^{r-1}_{m}\times X^{s-1}_{m}$, and let $Z$ be the
blow-up of $X^{r}_{m}\times X^{s}_{m}$ along $Y$. Shioda and Katsura show
that the rational map
\[
  (x,y)\mapsto\bigl[x_{0}y_{s+1}:\dots:x_{r}y_{s+1}:\varepsilon
  x_{r+1}y_{0}:\dots:\varepsilon x_{r+1}y_{s}\bigr]
\]
extends to a surjective morphism $\psi\colon Z\to X^{r+s}_{m}$
\cite[Lemma~1.2]{SK79}.\footnote{The image lies in $X^{r+s}_{m}$, because
$y_{s+1}^{m}\sum_{i\le r}x_{i}^{m}+\varepsilon^{m}x_{r+1}^{m}\sum_{j\le s}
y_{j}^{m}=-y_{s+1}^{m}x_{r+1}^{m}+x_{r+1}^{m}y_{s+1}^{m}=0$. All coordinates
vanish exactly on $Y$, and on the chart of the blow-up where
$y_{s+1}=t\,x_{r+1}$ the map is $[tx_{0}:\dots:tx_{r}:\varepsilon y_{0}:\dots:
\varepsilon y_{s}]$, where $y_{0},\dots,y_{s}$ do not all vanish.} Let
$E\subseteq Z$ be the exceptional divisor, $j\colon E\to Z$ the inclusion and
$\varpi\colon E\to Y$ the projection. The map
\[
  f=\psi_{*}j_{*}\varpi^{*}\colon H^{r+s-2}(Y)\to H^{r+s}(X^{r+s}_{m})
\]
carries $V(a')\otimes V(b')$ onto $V(a'*b')$ for $a'\in\mathfrak
A^{r-1}_{m}$ and $b'\in\mathfrak A^{s-1}_{m}$, where $a'*b'$ is the
juxtaposed character \cite[Theorem~II]{Shi79}, \cite[\S1]{Aok87}. As a
composite of pull-backs and push-forwards along morphisms, it carries
algebraic classes to algebraic classes.

\subsection{Balanced multisets}
In the rest of the paper $M\ge2$ is an integer, and we work with multisets
of elements of $\ZZ/M$; $A+B$ is the sum of multisets, $A-B$ the difference
when $B$ is contained in $A$, $tA$ and $-A$ the images under multiplication.

\begin{definition}\label{def:balanced}
Let $A$ be a multiset of nonzero elements of $\ZZ/M$ with sum $0$. It is
\emph{balanced} if $h_{A}(t)=\sum_{a\in A}\langle ta\rangle$ is the same
integer for every $t\in(\ZZ/M)^{\times}$. A \emph{pair} is a submultiset
$\{a,-a\}$; when $M$ is even, $\{M/2,M/2\}$ is a pair. A multiset is
\emph{split} if it is the sum of two multisets of three elements, each with
sum $0$.
\end{definition}

Comparing $t$ with $-t$ shows that a balanced multiset of $n$ elements has
$h_{A}=n/2$, so $n$ is even. Ordering the elements of a balanced multiset of
$n\ge2$ elements gives a character $\alpha\in\mathfrak A^{n-2}_{M}$, and by
\Cref{prop:fermatchars}(ii) the line $V(\alpha)$ consists of Hodge classes
(for $n=2$, $X^{0}_{M}$ is a set of $M$ points). Permuting the coordinates
permutes these lines, so the following notion does not depend on the order:
$A$ is \emph{algebraic} if $V(\alpha)\subseteq\Alg^{(n-2)/2}(X^{n-2}_{M})
\otimes\CC$. Conversely every Hodge class of $X^{4}_{M}$ is, over $\CC$, a
sum of a multiple of the square of the hyperplane class and of vectors of
lines $V(\alpha)$ with $\alpha$ balanced. Since $\Alg^{2}\subseteq\Hdg^{2}$
are rational subspaces, this gives:

\begin{lemma}\label{lem:reducetosextuples}
The Hodge conjecture holds for $X^{4}_{M}$ if and only if every balanced
multiset of six nonzero elements of $\ZZ/M$ is algebraic.
\end{lemma}

A balanced multiset of six nonzero elements is a \emph{sextuple at level
$M$}.

For an odd prime $\ell\mid M$ and $x\in\ZZ/M$ put
$\Phi_{\ell}(x)=\{x+jM/\ell:\ 0\le j<\ell\}$, the solutions $y$ of
$\ell y=\ell x$, and
\[
  \sigma_{\ell,x}=\Phi_{\ell}(x)+\{-\ell x\}.
\]
For even $M$ and $x\in\ZZ/M$ with $2x\ne0$ put
\[
  T(x)=\{x,\ x+M/2,\ -2x,\ M/2\}.
\]
These are the \emph{standard multisets} of Aoki \cite[\S1]{Aok87}.

\begin{lemma}\label{lem:algtools}
Let $A$ and $B$ be balanced multisets of nonzero elements of $\ZZ/M$.
\begin{enumerate}[label=\textup{(\alph*)}]
\item If $|A|\le4$, then $A$ is algebraic.
\item If $A$ and $B$ are algebraic, so is $A+B$. If $P$ is a sum of pairs and
$A+P$ is algebraic, then $A$ is algebraic.
\item For a unit $t$, $tA$ is balanced, and it is algebraic if and only if
$A$ is.
\item Let $g\mid M$, $M'=M/g$, let $A'$ be a multiset of nonzero elements of
$\ZZ/M'$ with sum $0$, and let $gA'\subseteq\ZZ/M$ be its image under
$y\mapsto gy$. Then $gA'$ is balanced if and only if $A'$ is, and in that
case $gA'$ is algebraic if and only if $A'$ is.
\item If $\ell x\ne0$, then $\sigma_{\ell,x}$ is balanced and algebraic. If
$M$ is even and $2x\ne0$, then $T(x)$ is balanced and algebraic.
\item A split sextuple is algebraic.
\end{enumerate}
\end{lemma}

\begin{proof}
(a) For $|A|=2$, $A=\{a,-a\}$ and $H^{0}(X^{0}_{M})$ is spanned by classes of
points. For $|A|=4$, $V(\alpha)$ consists of Hodge classes of degree two on
the surface $X^{2}_{M}$, which are algebraic by \Cref{thm:lef11}.

(b) The line of $A+B$ is the image of $V(\alpha)\otimes V(\beta)$ under the
map $f$ of the inductive structure, which preserves algebraic classes; this
is \cite[Theorem~1-4]{Aok87}, which also gives the converse for a pair
$\{a,-a\}$, whose line on $X^{0}_{M}$ is spanned by a class of points.

(c) The automorphism $\zeta\mapsto\zeta^{t}$ of $\QQ(\mu_{M})$, acting on
$H^{*}(X^{n-2}_{M},\QQ)\otimes\QQ(\mu_{M})$ through the second factor,
carries $V(\alpha)$, which is defined over $\QQ(\mu_{M})$, onto $V(t\alpha)$,
and preserves $\Alg\otimes\QQ(\mu_{M})$.

(d) Every unit modulo $M'$ lifts to a unit modulo $M$, and
$\langle tgy\rangle=\langle t'y\rangle$ for $y\in\ZZ/M'$ and a unit $t$ with
image $t'$, the first bracket taken modulo $M$ and the second modulo $M'$;
this gives the first assertion. Let $\pi\colon X^{n-2}_{M}\to X^{n-2}_{M'}$
be the finite morphism $[x_{i}]\mapsto[x_{i}^{g}]$, of degree
$g^{n-1}$. It is equivariant for the homomorphism $G\to G'$,
$\zeta\mapsto\zeta^{g}$, so for $\xi\in V(\alpha')$ and $\zeta\in G$,
$\zeta^{*}\pi^{*}\xi=\pi^{*}(\zeta^{g})^{*}\xi=\alpha'(\zeta^{g})\pi^{*}\xi$,
and $\pi^{*}V(\alpha')\subseteq V(g\alpha')$. Since $\pi_{*}\pi^{*}$ is
multiplication by $\deg\pi$, $\pi^{*}$ is injective, and as both are lines,
$\pi^{*}V(\alpha')=V(g\alpha')$ and $V(\alpha')=\pi_{*}V(g\alpha')$. Both
$\pi^{*}$ and $\pi_{*}$ preserve algebraic classes.

(e) The elements of $\Phi_{\ell}(x)$ are nonzero since $\ell x\ne0$, and the
sum of $\sigma_{\ell,x}$ is $\ell x+\binom{\ell}{2}M/\ell-\ell x=0$, as $\ell$
is odd. For a unit $t$, $t\Phi_{\ell}(x)=\Phi_{\ell}(tx)$, and writing
$\overline{tx}=qM/\ell+r$ with $0\le r<M/\ell$,
\[
  \sum_{y\in\Phi_{\ell}(tx)}\bar y=\sum_{j=0}^{\ell-1}\Bigl(\bigl((q+j)\bmod
  \ell\bigr)\frac{M}{\ell}+r\Bigr)=\frac{(\ell-1)M}{2}+\overline{\ell tx};
\]
adding $\overline{-\ell tx}=M-\overline{\ell tx}$ gives $(\ell+1)M/2$, so
$\sigma_{\ell,x}$ is balanced. For $\ell=3$ it is algebraic by (a). Let
$\ell\ge5$, $d=M/\ell$, and let $i$ be the representative of $x$ in
$\{1,\dots,d-1\}$ modulo $d$ (it is not $0$, as $\ell x\ne0$); then
$\Phi_{\ell}(x)=\Phi_{\ell}(i)$ and $\ell x=\ell i$. Put $g=\gcd(i,d)$,
$M'=M/g$, $d'=d/g$, $i'=i/g$. Then $\sigma_{\ell,x}=g\,\sigma_{\ell,i'}$ with
$\sigma_{\ell,i'}$ formed at level $M'$, and $\gcd(i',d')=1$, $d'\ge2$. If
$d'\ge3$, Aoki constructs a subvariety of $X^{\ell-1}_{M'}$ of codimension
$(\ell-1)/2$ that represents $\sigma_{\ell,i'}$ \cite[Theorem~2-1]{Aok87}, so
$\sigma_{\ell,i'}$ is algebraic. If $d'=2$, then $i'=1$, $M'=2\ell$, and
$\sigma_{\ell,1}=\{1,3,\dots,2\ell-1\}+\{\ell\}$ is the sum of the pairs
$\{y,2\ell-y\}$, $y<\ell$ odd, and $\{\ell,\ell\}$, algebraic by (a) and (b).
In both cases (d) gives the claim.

For $T(x)$: let $t$ be a unit, which is odd, and $y=\overline{tx}$; then
$y\ne M/2$ and $tM/2=M/2$. The representatives of the elements of $tT(x)$
are $y,\ y+M/2,\ M-2y,\ M/2$ if $y<M/2$, and $y,\ y-M/2,\ 2M-2y,\ M/2$ if
$y>M/2$; both add up to $2M$. So $T(x)$ is balanced, and it is algebraic by
(a).

(f) Let $A=A_{1}+A_{2}$ with $|A_{i}|=3$ and sums $0$. Their elements are
nonzero, and the entries of $A_{1}$ and $A_{2}$ are characters $a'$ and $b'$
of the Fermat curve $X^{1}_{M}$, with $\langle ta'\rangle,\langle
tb'\rangle\in\{1,2\}$ and $\langle ta'\rangle+\langle tb'\rangle=3$ for every
unit $t$. So $V(ta')\otimes V(tb')$ has type $(1,1)$ in
$H^{2}(X^{1}_{M}\times X^{1}_{M})$, and the sum $L$ of these lines over the
units is stable under $\operatorname{Gal}(\QQ(\mu_{M})/\QQ)$, so
$L=L_{\QQ}\otimes\CC$ for a space $L_{\QQ}$ of Hodge classes of degree two on
the surface $X^{1}_{M}\times X^{1}_{M}$. They are algebraic by
\Cref{thm:lef11}, and $f$ with $r=s=2$ carries $V(a')\otimes V(b')$ onto the
line of $A$.
\end{proof}

\subsection{Moves}
The proofs below replace a sextuple by another multiset in a way that keeps
track of algebraicity.

\begin{definition}\label{def:move}
Let $A$ be a balanced multiset of nonzero elements, $\sigma$ an algebraic
balanced multiset, and $S\subseteq A$ a submultiset that is also contained in
$\sigma$. The multiset
\[
  A''=(A-S)+\bigl(-(\sigma-S)\bigr)
\]
is the result of the \emph{move} $(\sigma,S)$. We usually take $\sigma$
standard, $S$ consisting of the elements of $A$ that lie in $\sigma$, and
we remove from $A''$ every pair it contains.
\end{definition}

\begin{lemma}\label{lem:move}
In \Cref{def:move}, $A''$ is a balanced multiset of nonzero elements and
$A''+\sigma=A+P$ for a sum of pairs $P$. If $A''$ is algebraic, so is $A$.
The same holds for a chain of moves.
\end{lemma}

\begin{proof}
With $P=(\sigma-S)+(-(\sigma-S))$ we have $A''+\sigma=(A-S)+S+(\sigma-S)+
(-(\sigma-S))=A+P$. The sum and the function $h$ are additive, and $A$, $P$,
$\sigma$ are balanced; so $A''$ is balanced. If $A''$ is algebraic, then
$A+P=A''+\sigma$ is algebraic by \Cref{lem:algtools}(b), and so is $A$.
\end{proof}

The \emph{profile} $o(A)$ of a sextuple is the sequence of the additive
orders of its elements in nonincreasing order, and profiles are compared
lexicographically. For two sextuples $A$ and $A''$, $o(A'')<o(A)$ exactly when,
at the largest integer $f$ at which $A$ and $A''$ have different numbers of
elements of order $f$, $A$ has more. A chain of moves that turns a sextuple
$A$ into a multiset $A^{\flat}$ with at most four elements is \emph{direct};
it is \emph{lowering} if $A^{\flat}$ is a sextuple with $o(A^{\flat})<o(A)$;
it \emph{succeeds} if it is direct or lowering, or $A^{\flat}$ is a split
sextuple, or a sextuple that lies in a proper subgroup of $\ZZ/M$.

\begin{proposition}[Induction]\label{prop:induction}
Let $M\ge2$, and suppose that the Hodge conjecture holds for $X^{4}_{M'}$
for every proper divisor $M'$ of $M$. Suppose that every sextuple at level
$M$ that contains no pair, is not split and generates $\ZZ/M$ has a chain of
moves that succeeds, or is algebraic for another reason. Then the Hodge
conjecture holds for $X^{4}_{M}$.
\end{proposition}

\begin{proof}
By \Cref{lem:reducetosextuples} we show that every sextuple $A$ at level $M$
is algebraic, by induction on the profile, which takes finitely many values.
If $A$ contains a pair, then removing it leaves a balanced multiset of four
elements, and $A$ is algebraic by \Cref{lem:algtools}(a),(b). If $A$ is split,
\Cref{lem:algtools}(f) applies. If $A$ lies in the subgroup $g\ZZ/M$, $g>1$,
then $A=gA'$ for a sextuple $A'$ at level $M/g$, which is algebraic by
\Cref{lem:reducetosextuples} and the hypothesis, and so is $A$ by
\Cref{lem:algtools}(d). Otherwise $A$ has a chain of moves that succeeds, or
is algebraic by assumption. In the first case the last multiset $A^{\flat}$
is algebraic by \Cref{lem:algtools}(a), by induction on the profile, or by
the cases already treated, and $A$ is algebraic by \Cref{lem:move}.
\end{proof}
